\documentclass[12pt]{article}
\usepackage[utf8]{inputenc}
\usepackage[T1]{fontenc}
\usepackage[margin=1in]{geometry}
\usepackage{setspace}
\usepackage{xcolor}
\usepackage[colorlinks=true, urlcolor=blue!60!black, linkcolor=blue!60!black]{hyperref}
\setlength{\parindent}{0em}
\setlength{\parskip}{0.6em}
\renewcommand{\baselinestretch}{1.0}
\begin{document}
\begin{center}
{\LARGE \textit{Printz v. United States}, 521 U.S. 898 (1997)\par}
\vspace{0.5em}
{\large Case Brief --- Two-Page Edition\par}
\end{center}

\section{Memory Jogger}
Congress may not conscript state or local executive officers to administer a federal program; the Brady Act's interim background-check mandate fell (\href{https://www.law.cornell.edu/supct/html/95-1478.ZO.html}{521 U.S.\ 898}).

\section{Detailed Case Facts}
Until a national instant-check system was ready, the Brady Act of 1993 required each local ``chief law enforcement officer'' (CLEO) to ``make a reasonable effort to ascertain within 5 business days'' whether a prospective handgun buyer could lawfully buy. Sheriffs Jay Printz of Montana and Richard Mack of Arizona sued, objecting to being pressed into federal service.

\section{Procedural History}
Both district courts held the mandate unconstitutional but severable; a divided Ninth Circuit reversed.

\section{Judicial Votes}
Reversed, 5--4. Scalia wrote for Rehnquist, O'Connor, Kennedy, and Thomas; O'Connor and Thomas concurred; Stevens (joined by Souter, Ginsburg, Breyer), Souter, and Breyer dissented.

\section{Holding}
Congress may neither direct the states to address problems nor command state officers to administer or enforce a federal regulatory program. The rule is categorical, with no balancing. The Court reserved purely ministerial reporting requirements.

\section{Analysis of Opinions}
Admitting no text addressed the question, Scalia found that early Congresses imposed duties on state judges but not executive officers, read \textit{Federalist} No.\ 27 as requiring only non-obstruction, and grounded the rule in dual sovereignty, political accountability, and the unitary executive; \textit{New York v.\ United States} controlled. O'Connor noted the Spending Clause alternative; Thomas questioned commerce power over intrastate sales. Stevens relied on the Supremacy Clause, founding expectations, and political safeguards; Souter on \textit{Federalist} No.\ 27; Breyer on European federal systems.

\section{Examples: Future Applications}
\textbf{Same-Side:} Congress may not force local police to honor federal immigration detainers. \textbf{Opposite-Side:} federal grants conditioned on voluntary background checks are valid. \textbf{Fence-Sitter:} a mandate that sheriffs report missing-children statistics, the question the Court reserved.

\section{Critique}
Formalists praise the defense of dual sovereignty; critics such as Prakash note the thin textual basis and ambiguous history, and Jackson and Caminker fault the bright line, though sanctuary jurisdictions now rely on it.

\section{Key Quotations}
\begin{quote}
``We held in New York that Congress cannot compel the States to enact or enforce a federal regulatory program. Today we hold that Congress cannot circumvent that prohibition by conscripting the State's officers directly. The Federal Government may neither issue directives requiring the States to address particular problems, nor command the States' officers, or those of their political subdivisions, to administer or enforce a federal regulatory program. It matters not whether policymaking is involved, and no case by case weighing of the burdens or benefits is necessary; such commands are fundamentally incompatible with our constitutional system of dual sovereignty.''
\end{quote}
\noindent --- Justice Scalia, Opinion of the Court, 521 U.S.\ at 935.

\end{document}
